% Two small multiples: code size and output tokens per call, per language.
% Data: results/haiku-v1/attempts.jsonl (code_lines of attempt 0; output_tokens of all calls).
\begin{subfigure}[t]{0.49\textwidth}
\centering
\begin{tikzpicture}
\begin{axis}[
  reportchart, xbar, width=0.93\linewidth, height=5.6cm, bar width=8pt,
  xmin=0, xmax=300, xtick={0,100,200,300},
  xlabel={Mediana liczby linii kodu (próba 0)},
  ytick={1,...,7}, ymin=0.4, ymax=7.6,
  yticklabels={Prolog, Scala, Haskell, Smalltalk, Java, Pascal, C},
  xmajorgrids=true, axis x line*=bottom, axis y line*=left,
  y axis line style={draw=none}, ytick style={draw=none},
]
  \addplot[fill=cat1, draw=none] coordinates {(248.5,7) (273.5,6) (184.5,5) (132,4) (127,3) (104.5,2) (160,1)};
\end{axis}
\end{tikzpicture}
\caption{Rozmiar rozwiązania.}
\label{fig:size-loc}
\end{subfigure}\hfill
\begin{subfigure}[t]{0.49\textwidth}
\centering
\begin{tikzpicture}
\begin{axis}[
  reportchart, xbar, width=0.93\linewidth, height=5.6cm, bar width=8pt,
  xmin=0, xmax=10, xtick={0,2,...,10},
  xlabel={Tokeny wyjściowe na wywołanie, tys.\ (z rozumowaniem)},
  ytick={1,...,7}, ymin=0.4, ymax=7.6,
  yticklabels={Prolog, Scala, Haskell, Smalltalk, Java, Pascal, C},
  xmajorgrids=true, axis x line*=bottom, axis y line*=left,
  y axis line style={draw=none}, ytick style={draw=none},
]
  \addplot[fill=cat1, draw=none] coordinates {(6.014,7) (9.320,6) (3.791,5) (7.166,4) (4.762,3) (4.459,2) (9.463,1)};
\end{axis}
\end{tikzpicture}
\caption{Nakład na generowanie.}
\label{fig:size-tokens}
\end{subfigure}
